\chapter{State of the Art}
\label{state-of-the-art}

% ""
% Overview of P2P file-sharing protocols (BitTorrent, Gnutella, etc.)

% Introdution to BitTorrent

\section{BitTorrent Protocol}
\label{BitTorrent}

Traditional HTTP downloads rely on centralized servers (CDN), creating a single point of failure where all operation cost falls on the hosting server. In contrast, BitTorrent, one of the most widely used P2P protocols, distributes this load across multiple peers in a swarm, significantly reducing the burden on individual participants.

BitTorrent achieves this distribution by splitting files into fixed-size chunks (typically 256KB to 4MB, depending on file size). For each chunk, a SHA1 is calculated to verify integrity. These hashes are stored in a .torrent metadata file, which peers use to verify that downloaded chunks are authentic and uncorrupted. Peers then exchange these chunks, propagating them accros the swarm until all peers possess complete copies.

%% .torrent structure %%
\subsection{The .torrent Metadata File}
\label{torrent_file}
Before peers can participate in a swarm, they must obtain the \textbf{.torrent} file, which serves as the blueprint for the entire download process. This file contains all necessary metadata encoded in Benconding format.

A typical .torrent file includes:

\begin{itemize}
  \item \textbf{announce:} The URL(s) of tracker(s) that coordinate the swarm
  \item \textbf{info:} A dictionary containing:
    \begin{itemize}
      \item \textbf{name:} The suggested filename or directory name
      \item \textbf{piece length:} The size of each chunk (typically 256KB-4MB)
      \item \textbf{pieces:} A concatenated string of SHA-1 hashes for all chunks
      \item \textbf{length/files:} File size (single file) or file list (multi-file torrents)
    \end{itemize}
  \item \textbf{nodes:} (Optional) Bootstrap nodes for DHT-based peer discovery
  \item \textbf{info hash:} A SHA-1 hash of the info dictionary, serving as
    the unique identifier for the torrent
\end{itemize}

Figure 1.1 shows examples of .torrent files using different peer discovery mechanisms.

The \textbf{.torrent} file is typically created by the initial seeder and distributed through external channels (websites, forumns, etc...). Also, the tracker URL and peer discovery method embbeded in this file determine how peers will interact with the network, each distinct security and privacy issues.

To understand how this distributed system operates, it is essencial to identity the different roles that peers assume within a BitTorrent swarm, each serving distint functions in file distribution process.

\subsection{Roles in the BitTorrent Ecosystem}
The BitTorrent network consists of several key participants:

\begin{itemize}
  \item \textbf{Seeders}: Peers that posses 100\% of the file and only upload chunks to others. At least one seeder must exists for a torrent to remain viable; without seeders, new peers cannot obtain a complete copy. Seeders are critical for swarm health and longevity.

  \item \textbf{Leechers (Downloaders):}: Peers with incomplete copies of the file (less then 100\%). They simultaneously download missing chunks whule uploading chunks they already verified. The term "leecher" sometimes carries negative connotations, refering to peers that download without adequately uploading in return, violating the reciprocity principle that BitTorrent relies upon.

  \item \textbf{Trackers:} Specialized servers that do not store or transfer file content. Instead, they maintain lists of active peers in each swarm and facilitate peer discovery by providing IP addresses of other participants. When a peer joins a swarm, it announces itself to the tracker, which responds with a list of other peers to connect to. However, this creates a privacy concern: trackers can log all IP addresses and associate them with specific content being downloaded.
\end{itemize}

While trackers have historically been the primary peer discovery mechanism, their centralized nature presents both operational vulnerabilities (single points of failure) and privacy risks (activity logging). This has led to the development of descentralized alternatives, discussed in <Section (peer discovery)>

\subsection{Peer Discovery Mechanisms}
\label{peer_discovery}
%%[Trackers, PEX, DHT]

\subsection{Data Exchange Strategies}
\label{data_exchange}
%% speak about Rarest First and Choking Algorithms

Once peers are connected, BitTorrent employs specific algorithms to
optimize distribution and enforce fairness.

\subsubsection{Rarest First Strategy}

Peers prioritize downloading the rarest pieces in the swarm first...
[1 paragraph explanation]
[Brief security note: helps ensure piece availability]

\subsubsection{Choking Algorithm}

To prevent freeloading and optimize bandwidth usage, BitTorrent uses a
tit-for-tat mechanism...
[1-2 paragraphs]
[Security relevance: trust/reputation, but vulnerable to strategic behavior]

%% -- example --

%% utilizam bson - um compressor especial no BitTorrent
%% 1. mostrar em json format vários casos, quando em:
%% 1.1. Single File
%% 1.2. Tracker
%% 1.3. DHT
%% ---
%% Merkle Trees, para calcular a integridade dos próprios hashs dos chunks
%% https://en.wikipedia.org/wiki/Torrent_file <- explain
%% https://www.bittorrent.org/beps/bep_0003.html <- another font, but it's not wikipedia

\subsection{Security and Privacy Issues in BitTorrent}
\label{issues_bittorrent}
%%[Lead naturally into your motivation for TEEs, QUIC, and Onion routing]

%% Existing attempts at secure or anonymous P2P.

%
% 1. dates, usage, speak how it's used nowadays, for what porpuses?, etc...
%%O BitTorrent é um dos protocolos P2P mais utilizados hoje em dia, permitindo distribuir conteúdos multimédia com grandes tamanhos facílmente, supostamente mais rápidos do que efetuar a transferência de um servidor único. Quando um fichiero é disponibliza por um link magnetico HTTP, normalmente este é tranferidode um único servidor, normalmente um CDN (grab label) próximo do utilizador, desta forma todo o custo da operação recai sobre um único host (servidor). Normalmente, existem várias técnicas para fazer com que a transferência destes contêudos seja eficiênte, normalmente estes vêm aconpanhados com um fichiro que armazena os metadados do ficheiro que utilizador está a tentar realizar transferência, que contêm informações sobre o ficheiro, como: nome, tamanho, tipo de ficheiro, etc... Além disso, para facilitar o download e aproveitar ao máximo a largura de banda do utilizador para o servidor, este ficheiro é repartido em vários chunks de um tamanho fixo, e de seguida são calculados para cada um os hashes e demarcados no seu metafile, depois quando um cliente pertende tranferir o ficheiro. começa por tranferir o seu metafile, para obter acesso q


% 2. Speak how it works, overview -> structure of .torrent all of it, kinda explain how it works too!
%   3.1 Speak about Trackers -> correlate with Tor, and TEE (Peer Discovery)
%   3.2 Speak about PEX ->  did not understand quite this, but some day i will (Explain more or less, and it's limitations, paper cite)
%   3.3 Speak about DHT -> speak about sybil attacks and eclipse attacks (explain basecly kademlia, cite it, and give examples of attacks and how to mitigate them at least my response to it)
% 4. Choking Algorithms
% 4. Rarest First
